\chapter{Introducción}
\label{ch:introduccion}

En un contexto económico global en el que algunos países enfrentan tasas de 
inflación elevadas, la adopción de monedas digitales como reserva de valor ha 
surgido como una alternativa atractiva para preservar el poder adquisitivo, 
intentando, de algún modo, evitar la erosión de valor que sufren las monedas locales.

Las criptomonedas ofrecen una opción descentralizada e internacional, 
proporcionando a los ciudadanos una alternativa que no está sujeta al control 
directo de los bancos centrales.

Debido a que la existencia de las criptomonedas no soluciona directamente el 
problema de la 
necesidad de que los ciudadanos sigan obligados a depositar su confianza para la 
financiación en bancos tradicionales y gobiernos centralizados, este proyecto surge 
para ofrecer una alternativa a esta necesidad.

Hasta ahora, las necesidades de gestión y confianza necesarias para un sistema de 
préstamos solo las podían cubrir grandes entidades centralizadas, pero gracias 
a la tecnología \gls{blockchain} y tal y como se explicará más adelante, 
estas características quedan a la disposición 
de cualquier ciudadano. 

El presente Proyecto Fin de Grado aborda el diseño, desarrollo y validación 
de un prototipo de gestión de préstamos basado en la red \gls{ethereum}. 
La propuesta utiliza \gls{eth} como activo para la concesión y devolución de los 
préstamos, y \glspl{smartcontract} para la gestión de estos.

Aunque el acceso a la financiación a través de este proyecto sea descentralizado,
la oferta de esta financiación sí queda centralizada con el objetivo de incentivar 
la producción de este sistema a una entidad competente.

La solución se organiza mediante un \gls{smartcontract} principal que administra los 
productos y los fondos disponibles y que genera \glspl{smartcontract} 
individuales que representan los préstamos contratados. 
En ellos se recoge la información necesaria para gestionar 
la deuda, las cuotas pendientes y las penalizaciones por impago. 
De este modo, la información del sistema que define la situación de los préstamos 
queda reflejada en la red \gls{ethereum} ofreciendo así las propiedades de 
veracidad e inmutabilidad necesarias para un sistema de este tipo.

Este proyecto analiza si un sistema de préstamos descentralizado basado en 
\gls{ethereum} puede ser accesible para los ciudadanos reduciendo la dependencia 
de financiación con modelos 
tradicionales.

El trabajo se centrará en la creación de un prototipo funcional capaz de 
comprobar la viabilidad técnica de esta propuesta haciendo uso de pruebas 
sobre los \glspl{smartcontract} de \gls{ethereum} y validando los flujos 
de uso de la aplicación.




\section{Objetivos}

En esta sección se detallan los objetivos de este proyecto, 
estos se dividen en dos objetivos principales  (\gls{op}) y 
distintos objetivos secundarios (\gls{os}) los cuales surgen en el desarrollo 
de los objetivos principales.

\newcounter{objetivoprincipal}
\newcounter{objetivosecundario}

\subsection{Objetivos principales}

\begin{itemize}
    \item \refstepcounter{objetivoprincipal}\label{op:1}\textbf{\gls{op} 1}: Desarrollo de un sistema de préstamos automatizado 
    sin una entidad centralizada controladora, de acceso libre y cumpliendo 
    los requisitos básicos de un proceso de préstamo, como pueden ser la seguridad, 
    historial de cuotas, seguimiento, gestión de morosidad y personalización 
    del usuario.
    
    \item \refstepcounter{objetivoprincipal}\label{op:2}\textbf{\gls{op} 2}: Desarrollar una interfaz unificada que permita a los usuarios, 
    con o sin conocimientos técnicos, interactuar con el sistema de manera fácil e 
    intuitiva. La interfaz garantizará que todas las funcionalidades sean accesibles 
    desde un solo punto, facilitando la experiencia del usuario.
\end{itemize}

\subsection{Objetivos secundarios}
\begin{itemize}
    \item \refstepcounter{objetivosecundario}\label{os:1}\textbf{\gls{os} 1}: Explorar la viabilidad de los sistemas de identificación 
    en aplicaciones \gls{defi}.

    \item \refstepcounter{objetivosecundario}\label{os:2}\textbf{\gls{os} 2}: Crear una solución segura y eficiente.  
    Almacenar únicamente la información estrictamente necesaria para el funcionamiento 
    correcto de la aplicación evitando así la exposición de datos sensibles y el procesamiento 
    innecesario.

    \item \refstepcounter{objetivosecundario}\label{os:3}\textbf{\gls{os} 3}: Optimizar la eficiencia de ejecución del código en
    la red blockchain para reducir la cantidad de \gls{gas} 
    necesario para el correcto funcionamiento del sistema.

    \item \refstepcounter{objetivosecundario}\label{os:4}\textbf{\gls{os} 4}: Crear mecanismos de control de  cumplimiento de contrato 
    y penalización automatizados.
    
    \item \refstepcounter{objetivosecundario}\label{os:5}\textbf{\gls{os} 5}: Explorar la flexibilidad y personalización de 
    contratos inteligentes en productos financieros. 
    Se investigarán maneras en las que el sistema sea capaz de ofrecer 
    distintos productos financieros en función de la información que el sistema tenga 
    del cliente como por ejemplo su situación financiera. 
    
    \item \refstepcounter{objetivosecundario}\label{os:6}\textbf{\gls{os} 6}: Optimizar el equilibrio entre la naturaleza pública y transparente
    de una red \gls{blockchain} y la privacidad necesaria para los usuarios. 
    
    \item \refstepcounter{objetivosecundario}\label{os:7}\textbf{\gls{os} 7}: Analizar cómo se puede gestionar la liquidez en sistemas 
    \gls{defi} y cómo se  puede coordinar la transferencia de fondos en un sistema 
    de préstamos en una red \gls{blockchain}.
\end{itemize}

\section{Motivación}
\label{s:motivacion}

Demostrar la viabilidad de un sistema de préstamos descentralizado mostrando 
cómo la 
tecnología \gls{blockchain} y los \glspl{smartcontract} de \gls{ethereum} 
pueden ofrecer una alternativa 
sólida y segura frente a los sistemas financieros tradicionales, 
eliminando intermediarios y 
reduciendo  costos y barreras de acceso.

Explorar y afrontar los desafíos técnicos específicos de desplegar un sistema 
\gls{defi} 
en la red \gls{ethereum}.

Alcanzar la seguridad y eficiencia en los gastos del sistema, suponiendo esto 
una de las claves
para conseguir la rentabilidad económica en el proyecto y pudiéndose convertir
en uno de los 
puntos débiles del sistema si no se realiza con detalle.

La motivación también radica en el valor social y económico que una solución 
como esta puede aportar a un mundo en el que la inclusión financiera sigue siendo una meta distante
para millones de personas. 


Enrique Dans, doctorado en sistemas de la información, explica en su blog que los
proyectos \gls{defi} responden a la creciente demanda de alternativas que no 
solo sean económicamente 
accesibles, sino también transparentes y seguras, proporcionando una herramienta 
que permita a 
los usuarios mantener una mayor autonomía sobre sus decisiones financieras en 
un entorno digital 
global.~\cite{[3]}

\section{Estructura de la memoria}

El código fuente de este proyecto se puede encontrar en el siguiente 
\href{https://github.com/AkilinoGit/akilinogit.github.io}{repositorio}~\cite{[25]}.
El código fuente con las pruebas unitarias de este proyecto se puede encontrar en el siguiente
\href{https://github.com/AkilinoGit/prestamosFactoryMochaTest.git}{repositorio}~\cite{[34]}.

La estructura del documento es la siguiente:
\begin{itemize}
    \item \textbf{Introducción:}
    \begin{itemize}
        \item Presentación del problema y contexto económico de los préstamos descentralizados.
        \item Motivación y relevancia del proyecto.
        \item Objetivos generales y específicos del sistema.
    \end{itemize}
    \item \textbf{Estado de la cuestión}:
        \begin{itemize}
        \item Introducción a Blockchain y su importancia en el sector financiero.
        \item Explicación de Ethereum y su ecosistema.
        \item Tecnologías clave (Web3, Solidity, Remix IDE, Node.js, Truffle, etc.).
    \end{itemize}
    \item \textbf{Diseño de la Arquitectura:}
        \begin{itemize}
        \item Solución propuesta: Cómo se estructura el sistema de préstamos descentralizado.
        \item Suposiciones y limitaciones del proyecto.
        \item Actores del sistema y su interacción (Usuarios, Propietarios, Contratos).
        \item Casos de uso principales y pantallas del sistema.
        \item Diagramas UML, diagramas de secuencia y diagramas entidad-relación.
    \end{itemize}
    \item \textbf{Diseño detallado:}
            \begin{itemize}
        \item Explicación a nivel técnico de cada entidad del sistema.
        \item Smart Contracts implicados y su funcionamiento.
        \item Explicación detallada de la lógica del algoritmo principal.
    \end{itemize}
    \item \textbf{Conclusiones y trabajo futuro}
            \begin{itemize}
        \item Evaluación de los logros alcanzados en el proyecto.
        \item Posibles mejoras y nuevas funcionalidades a explorar.
    \end{itemize}
    
\end{itemize}
